\mxfront{expertes-front}
\mxcover{MATHÉMATIQUES EXPERTES}{Terminale générale · option}{Livret de cours et d'entraînement}{Nombres complexes · Arithmétique · Graphes et matrices}{2026-2027}
\frontmatter
\chapter*{Utiliser ce livret}
\addcontentsline{toc}{chapter}{Utiliser ce livret}
Ce livret couvre l'option mathématiques expertes de terminale générale. Il est construit pour être travaillé en parallèle de l'enseignement de spécialité. Chaque chapitre suit le même ordre.
\begin{enumerate}
\item \textbf{Automatismes} : quelques questions courtes pour réactiver les prérequis, à traiter sans calculatrice en début de séance.
\item \textbf{Activité d'introduction} : une recherche à mener avant le cours. On calcule, on observe, on formule une conjecture ; la preuve vient ensuite.
\item \textbf{Notes de cours} : définitions, propriétés, démonstrations et exemples. C'est la trace écrite à apprendre.
\item \textbf{Exercices} : classés par difficulté croissante. Les premiers vérifient une technique, les derniers demandent une rédaction complète.
\item \textbf{Problème} : une étude plus longue qui prolonge le chapitre.
\end{enumerate}
\begin{keypoint}[Expérimenter n'est pas démontrer]
Un calcul sur quelques exemples, une figure ou un tableau de valeurs donnent une conjecture. Seule une démonstration établit qu'elle est vraie dans tous les cas. Chaque activité se termine en le rappelant, comme chaque figure manipulable des diaporamas.
\end{keypoint}
\textbf{Programme suivi.} Programme de l'enseignement optionnel de mathématiques expertes de terminale générale, arrêté du 19 juillet 2019, BO spécial n\textsuperscript{o} 8 du 25 juillet 2019 (NOR MENE1921264A). Les démonstrations signalées comme exigibles par ce programme sont toutes rédigées dans les notes de cours. Ce qui dépasse le programme est placé dans un cadre « Hors programme ».

\textbf{Organisation.} Les chapitres sont regroupés par thème : nombres complexes (chapitres 1 à 5), arithmétique (chapitres 6 à 8), graphes et matrices (chapitres 9 à 11). En classe, il est conseillé d'alterner les thèmes, par exemple un chapitre de nombres complexes puis un chapitre d'arithmétique.
\ifmxteacher
\mxnewpage
\chapter*{Guide de préparation}
\addcontentsline{toc}{chapter}{Guide de préparation}
\begin{teachernote}[Contenu et éditions]
L'option de classe \texttt{edition=teacher} produit l'édition professeur, avec les corrigés des automatismes, de l'activité, des exercices et du problème, ainsi que des repères de conduite. L'option \texttt{edition=student} produit l'édition élève, sans corrigés. Les sources contiennent les deux éditions.
\end{teachernote}
\paragraph{Progression proposée.} Pour trois heures hebdomadaires, une alternance possible est : C1, C6, C2, C7, C3, C8, C4, C5, puis C9, C10, C11. Les chapitres de graphes et matrices utilisent les probabilités conditionnelles de spécialité ; les placer après leur étude en spécialité. Le chapitre 3 utilise le produit scalaire de première ; l'exercice de linéarisation intégrée du chapitre 3 suppose le calcul intégral de spécialité, il peut être reporté.
\paragraph{Durées.} Les durées des activités sont des estimations de travail effectif, non mesurées auprès d'élèves. Prévoir une lecture individuelle, une recherche par binômes, puis une mise en commun avant la trace écrite.
\paragraph{Démonstrations exigibles.} Conjugué d'un produit, d'un inverse, d'une puissance ; formule du binôme ; $\abs{z}^2=z\overline{z}$ ; module d'un produit et d'une puissance ; une formule d'addition ; factorisations de $z^n-a^n$ et de $P(z)$ par $z-a$ ; nombre de racines au plus égal au degré ; détermination de $\mathbb{U}_n$ ; écriture du PGCD sous la forme $au+bv$ ; théorème de Gauss ; infinité des nombres premiers ; nombre de chemins de longueur $n$ et probabilité de transition en $n$ étapes. Elles figurent dans les notes de cours des chapitres correspondants.
\paragraph{Problèmes.} Les problèmes de fin de chapitre reprennent des « problèmes possibles » cités par le programme. Ils ne sont pas exigibles : ils servent au travail en autonomie, au devoir à la maison ou à l'approfondissement.
\fi
\tableofcontents
\mainmatter
